%==========================================================
%  CHAPITRE 3 — ANALYSE ET CONCEPTION
%==========================================================
\chapter{Analyse et conception}
\label{chap:conception}

%──────────────────────────────────────────────────────────
\section{Méthodologie de développement}
\label{sec:cpt:methodologie}
%──────────────────────────────────────────────────────────

\todo{Décrire la méthodologie adoptée (Scrum, Kanban, Cycle en V,
RUP, etc.) avec ses artefacts, ses rôles et ses rituels.}

\subsection{Processus Scrum}

\begin{center}
% Schéma simplifié du cycle Scrum
\begin{tikzpicture}[
  node distance=1.5cm,
  box/.style={
    rectangle, rounded corners=5pt,
    fill=ChapBlue!12, draw=ChapBlue,
    font=\sffamily\footnotesize\bfseries, text=ChapBlue,
    minimum width=2.8cm, minimum height=0.85cm,
    align=center
  },
  arrow/.style={-{Stealth[length=5pt]}, thick, draw=AccentBlue}
]
  \node[box, fill=ChapBlue, text=white] (PB)  {Product\\Backlog};
  \node[box, right=of PB]               (SP)  {Sprint\\Planning};
  \node[box, right=of SP]               (SB)  {Sprint\\Backlog};
  \node[box, right=of SB]               (DEV) {Sprint\\(2 sem.)};
  \node[box, right=of DEV]              (REV) {Sprint\\Review};
  \node[box, below=of DEV]              (RET) {Rétrospective};

  \draw[arrow] (PB) -- (SP);
  \draw[arrow] (SP) -- (SB);
  \draw[arrow] (SB) -- (DEV);
  \draw[arrow] (DEV) -- (REV);
  \draw[arrow] (REV) -- ++(0,-1.2) -| (SP);
  \draw[arrow] (DEV) -- (RET);
\end{tikzpicture}
\end{center}

\vspace{0.5em}

\begin{center}
\begin{tabular}{>{\sffamily\bfseries}l l p{7cm}}
  \toprule
  \textbf{Artéfact / Rôle} & \textbf{Responsable} & \textbf{Description} \\
  \midrule
  Product Backlog & \todo{Product Owner} & Liste priorisée de toutes les user stories \\
  Sprint Backlog  & Équipe              & Sélection de stories à réaliser dans le sprint \\
  Daily Standup   & Équipe              & Réunion quotidienne de 15 min \\
  Sprint Review   & Équipe + Stakeholders & Démo de l'incrément livré \\
  Rétrospective   & Équipe + Scrum Master & Amélioration continue du processus \\
  \bottomrule
\end{tabular}
\end{center}

%──────────────────────────────────────────────────────────
\section{Diagrammes UML}
\label{sec:cpt:uml}
%──────────────────────────────────────────────────────────

\subsection{Diagramme de cas d'utilisation}

\placeholder{Diagramme de cas d'utilisation global\\
(à insérer dans figures/usecase.png ou à générer avec PlantUML)}

La figure~\ref{fig:usecase} présente les interactions entre les
différents acteurs et le système.

\begin{figure}[H]
  \centering
  \placeholder{Diagramme de cas d'utilisation — figures/usecase.png}
  % \includegraphics[width=0.85\textwidth]{figures/usecase}
  \caption{Diagramme de cas d'utilisation global}
  \label{fig:usecase}
\end{figure}

\todo{Commenter chaque groupe de cas d'utilisation en précisant
l'acteur concerné et le déclencheur de l'interaction.}

\subsection{Diagramme de classes}

\begin{figure}[H]
  \centering
  \placeholder{Diagramme de classes — figures/class\_diagram.png}
  % \includegraphics[width=0.9\textwidth]{figures/class_diagram}
  \caption{Diagramme de classes du domaine}
  \label{fig:classes}
\end{figure}

\todo{Décrire les principales entités (classes), leurs attributs,
leurs méthodes et leurs relations (association, agrégation,
composition, héritage).}

\subsection{Diagramme de séquence — Scénario principal}

\begin{figure}[H]
  \centering
  \placeholder{Diagramme de séquence — scénario principal\\
               (figures/sequence\_main.png)}
  % \includegraphics[width=0.9\textwidth]{figures/sequence_main}
  \caption{Diagramme de séquence : \todo{Nom du scénario principal}}
  \label{fig:sequence}
\end{figure}

\todo{Expliquer le déroulement du scénario illustré.}

\subsection{Diagramme d'activité}

\begin{figure}[H]
  \centering
  \placeholder{Diagramme d'activité — figures/activity.png}
  % \includegraphics[width=0.75\textwidth]{figures/activity}
  \caption{Diagramme d'activité : \todo{Nom du processus métier}}
  \label{fig:activity}
\end{figure}

%──────────────────────────────────────────────────────────
\section{Architecture technique}
\label{sec:cpt:architecture}
%──────────────────────────────────────────────────────────

\subsection{Architecture globale}

\todo{Décrire l'architecture retenue (monolithique, microservices,
MVC, hexagonale, événementielle…) et justifier ce choix.}

\begin{figure}[H]
  \centering
  \begin{tikzpicture}[
    node distance=1.4cm and 2.2cm,
    layer/.style={
      rectangle, rounded corners=6pt,
      draw=#1, fill=#1!10,
      font=\sffamily\small\bfseries, text=#1,
      minimum width=4cm, minimum height=1cm,
      align=center
    },
    arrow/.style={-{Stealth[length=5pt]}, thick, draw=SecGrey}
  ]
    % Couches
    \node[layer=AccentBlue]               (frontend)  {Frontend\\(\todo{React / Angular})};
    \node[layer=ChapBlue, below=of frontend]  (backend)
                                                 {Backend API\\(\todo{Spring Boot / FastAPI})};
    \node[layer=SecGrey,  below=of backend]   (db)
                                                 {Base de données\\(\todo{PostgreSQL / MongoDB})};
    % Côté
    \node[layer=WarningOrange, right=3cm of backend] (ext)
                                                 {Services externes\\(\todo{Email, OAuth, S3})};
    % Flèches
    \draw[arrow] (frontend) -- node[right,font=\tiny\sffamily]{REST/HTTP} (backend);
    \draw[arrow] (backend)  -- node[right,font=\tiny\sffamily]{JDBC/ORM}  (db);
    \draw[arrow] (backend)  -- node[above,font=\tiny\sffamily]{API calls} (ext);
  \end{tikzpicture}
  \caption{Architecture en couches du système}
  \label{fig:architecture}
\end{figure}

\subsection{Architecture de base de données}

\todo{Présenter le schéma de la base de données (MCD / MLD) ou décrire
les collections si NoSQL.}

\begin{figure}[H]
  \centering
  \placeholder{Schéma entité-relation / MLD — figures/db\_schema.png}
  % \includegraphics[width=0.9\textwidth]{figures/db_schema}
  \caption{Schéma de la base de données}
  \label{fig:db}
\end{figure}

\subsection{Description des composants principaux}

\begin{description}[leftmargin=0pt]
  \item[\textcolor{ChapBlue}{\sffamily\bfseries Module A — \todo{Nom}}]
    \todo{Décrire le rôle, les responsabilités et les interfaces
    exposées par ce module.}

  \item[\textcolor{ChapBlue}{\sffamily\bfseries Module B — \todo{Nom}}]
    \todo{Décrire le rôle, les responsabilités et les interfaces
    exposées par ce module.}

  \item[\textcolor{ChapBlue}{\sffamily\bfseries Module C — \todo{Nom}}]
    \todo{Décrire le rôle, les responsabilités et les interfaces
    exposées par ce module.}
\end{description}

%──────────────────────────────────────────────────────────
\section{Modélisation des données}
\label{sec:cpt:donnees}
%──────────────────────────────────────────────────────────

\todo{Si applicable, décrire les principales entités du domaine,
leur modélisation en base de données et les règles de gestion associées.}

\begin{center}
\begin{tabular}{>{\sffamily\bfseries}l >{\ttfamily}l l l}
  \toprule
  \textbf{Entité} & \textbf{Table} & \textbf{Clé primaire} & \textbf{Description} \\
  \midrule
  \todo{Utilisateur} & \todo{users}    & id (UUID)  & \todo{Description} \\
  \todo{Commande}    & \todo{orders}   & id (SERIAL)& \todo{Description} \\
  \todo{Produit}     & \todo{products} & id (SERIAL)& \todo{Description} \\
  \todo{Entité 4}    & \todo{table\_4} & id         & \todo{Description} \\
  \bottomrule
\end{tabular}
\end{center}

\clearpage
